% 第2章
\part{壁面走行システムの設計と製作}
\label{design}

\section{機体コンセプトと足回り機構}

\subsection{機体諸元}
機体の主要諸元は以下の通りである．
  \begin{table}[h]
      \centering
      \caption{Specifications of the system} % ここで自動的に Table 1 と表示されます
      \label{tab:spec}
      \begin{tabular}{ll}
      \hline
      Size           & 400 $\times$ 300 $\times$ 200 (\,mm) \\ \hline
      Weight         & 8.18 (\,kg) \\ \hline
      Motor          & DC Motor (MS-555VC) $\times$ 2 \\ \hline
      Power source   & \begin{tabular}[c]{@{}l@{}}Li-Poly Battery\\ (14.8V, 4000mAh, 100C)\end{tabular} \\ \hline
      Wheels         & \begin{tabular}[c]{@{}l@{}}Urethane Wheels\\ (ROEKGA100-20-50-T10) $\times$ 2\end{tabular} \\ \hline
      Controller     & \begin{tabular}[c]{@{}l@{}}Wireless Controller\\ (PS4 DualShock 4)\end{tabular} \\ \hline
      \end{tabular}
 \end{table}
  

\newpage

\subsection{三点支持構造による接地安定性の確保}
壁面走行ロボットにおいて，
全車輪が常に壁面に接地しトルクを伝達することは，
安全な走行を実現する上で不可欠である．
一般的な四輪構造では，壁面に微細な不陸（凹凸）が存在する場合，
対角の車輪が浮き上がる「がたつき」が生じやすい．
本研究では，これを解決するために2つの駆動輪と1つのフリーローラからなる三点支持構造を採用した．
三点支持は幾何学的に常に一つの平面を決定する静定構造であるため，
複雑なサスペンション機構を用いることなく常に全車輪が接地し，
安定した接地圧を確保することが可能である．

\begin{figure}[H] 
  \centering 
  \includegraphics[width=0.6\textwidth]{santen.png} 
  \vspace{-2mm} 
  \caption{Three-point support mechanism using industrial urethane rollers} \label{fig:result} 
  \end{figure}

\newpage


\subsection{産業用ウレタンローラの採用と選定理由}
駆動輪およびフリーローラには，
産業用規格品のウレタンローラを採用した．
これは，以下の実務的な観点に基づいた選定である．
\begin{itemize}
    \item \textbf{高いグリップ力：} 鋼材壁面に対して高い摩擦係数を持つウレタンは，剥離反力に抗して機体を前進させるのに適している．
    \item \textbf{耐摩耗性とメンテナンス性：} 実橋梁の現場は錆や研削材が飛散する過酷な環境である．摩耗した際に，汎用品として安価かつ迅速に交換できるメンテナンス性の高さ（ミスミ等での入手性）を重視した．
    \item \textbf{衝撃吸収：} 適度な弾性により，段差乗り越え時の衝撃を和らげる効果がある．
\end{itemize}

\newpage

\subsection{ベルト駆動方式の採用}
動力伝達にはタイミングベルトを用いた間接駆動方式を採用した．
モータ軸と車輪軸を分離することで，機体内部のスペースを有効活用し，
重心を低く抑える設計とした．
また，直接駆動に比べてモータへの衝撃伝達を抑えられるメリットがある．

\newpage

\section{フレーム設計と実務的知見の反映}
\subsection{アルミプロファイルによる高剛性フレーム}
機体メインフレームには，軽量かつ比強度の高いアルミプロファイルを採用した．本機体の設計にあたっては，実用的な産業用ロボットの設計指針に基づき，過酷な作業環境下でも安定した性能を維持するための以下の工夫を施している．

1. 磁気吸着荷重に対する剛性の最適化 本機体は強力なネオジム磁石により常時壁面方向へ引きつけられており，フレームには常に大きな圧縮荷重と曲げモーメントが作用する．これに対し，接合部には肉厚の大型直角ブラケットを配置することで，フレームの弾性変形を抑制し，幾何学的な直角度を徹底して確保した．これにより，駆動軸の歪みを防ぎ，安定した直進走行性能を維持している．

2. 現場におけるモジュール性とメンテナンス性 実橋梁での実験では，現場の状況に応じてセンサやスクレーパの取り付け位置を微調整する必要がある．アルミプロファイルの溝を利用したボルト締結構造を採用することで，溶接構造にはない「高い拡張性（モジュール性）」を確保した．これは，部品の破損時や機構の改良時において，現場での迅速な部品交換や位置調整を可能にし，実験の継続性を高めるための実務的な設計判断である．

\begin{figure}[H] 
  \centering 
  \includegraphics[width=0.6\textwidth]{ashimawari004.png} 
  \vspace{-2mm} 
  \caption{Aluminum profile frame structure and modular assembly} \label{fig:result} 
  \end{figure}
\newpage

\section{磁気吸着ユニットの設計}
\subsection{磁石配置の非対称化（12:8配置）}
本機体は合計20個のネオジム磁石（N52）を搭載している．
剥離反力は機体前方のスクレーパ部で発生し，
機体を引き剥がそうとするモーメントとして作用する．
このため，磁石配置を前方に12個，後方に8個という「非対称配置」とした．
これにより，前方の浮上を抑制する保持力を重点的に強化した．

\begin{figure}[H] 
  \centering 
  \includegraphics[width=0.6\textwidth]{haiti.png} 
  \vspace{-2mm} 
  \caption{Asymmetric layout of neodymium magnets (12 front, 8 back)} \label{fig:result} 
  \end{figure}
\newpage

\subsection{5mmの壁面距離（ギャップ）設定の根拠}
磁石面と壁面の距離は5mmに設定した．
磁石を壁面へ接近させるほど保持力は指数関数的に向上するが，
実橋梁では1mm程度の錆の隆起や塗膜の不陸が散見される．
1mm程度の距離に設定した場合，走行中に磁石ユニットがこれらに接触・激突し，
走行不能（スタック）に陥るリスクが高い．
そのため，現場環境の不確実性に対する「走破性」を優先し，
5mmという安全マージンを設定した．

\begin{figure}[H] 
  \centering 
  \includegraphics[width=0.6\textwidth]{gap.png} 
  \vspace{-2mm} 
  \caption{Side view of the robot showing the 5.0 mm clearance between the magnets and the wall surface} \label{fig:result} 
  \end{figure}
\newpage